\documentclass[a4paper,11pt]{article}
\usepackage[margin=2.3cm]{geometry}
\usepackage{amsmath,amssymb}
\usepackage{enumitem}
\usepackage{xcolor}
\usepackage[hidelinks]{hyperref}
\setlength{\parindent}{0pt}
\setlength{\parskip}{4pt}
\def\v#1{\mathbf{#1}}
\def\mf#1{\mathbb{#1}}
\newcommand{\Kv}{\mathcal K}
\newcommand{\MS}{\overline{\rm MS}}
\newcommand{\point}[1]{\par\medskip\noindent\textbf{#1}\par}
\newcommand{\referee}[1]{\begin{quote}\small\emph{Referee:} \emph{#1}\end{quote}}
\newcommand{\done}{\par\noindent\emph{What we did.}\ }
\newcommand{\result}{\par\noindent\emph{Result.}\ }
\newcommand{\where}{\par\noindent\emph{Where.}\ }
\newcommand{\open}{\par\noindent\emph{Open.}\ }

\title{Reply to Referee 3\\[4pt]\large Single inclusive gluon production at mid rapidity at next-to-leading order in the Color Glass Condensate}
\date{}
\begin{document}
\sloppy
\maketitle
\vspace{-1.5cm}

We thank the referee for a careful and constructive report. The referee re-derived a large part of the leading-order and leading-logarithmic structure and identified precisely the checks that were missing for the genuinely new NLO content. We have taken every point seriously. Several of them led to new results, two of which change statements of the manuscript in substance (the prescription of the vanishing energy denominators, Major~5, and the interpretation of the constant $\frac{67}{9}-\frac{\pi^2}{3}$, Major~3 and the error item on Section~7.6); others are answered by new checks or by rewording; some remain open, and we say so explicitly in a new subsection ``Open points'' (Section~8.1) instead of in red notes. Below we answer every bullet of the report in its order. For each item we state what we did, the result, where it is in the revised manuscript, and what remains open.

Section, equation and page numbers refer to the revised manuscript (118 pages instead of 100) unless they are marked ``old''. The section structure is unchanged; the old Sections~6.1--6.8 are now 6.2--6.9 (a new Section~6.1 contains the compact summary of the result), the old Sections~7.2 and 7.3 are now 7.3 and 7.4, and the old Sections~7.4--7.8 are now 7.6--7.10, with two new subsections (7.2, the dilute--dilute limit; 7.5, cubic terms on exact charge states). The old Eq.~(7.2) is now Eq.~(7.5).

\paragraph{Summary of the main changes.}
\begin{itemize}[leftmargin=1.5em,itemsep=1pt]
\item Compact summary of the complete result with pointers to all equations (Section~6.1, Eq.~(6.2) and table).
\item Dilute--dilute limit (new Section~7.2): the real emission is \emph{exactly} the quasi-multi-Regge two-gluon production cross section of Fadin--Lipatov in the form of Bartels--Sabio Vera--Schwennsen; for a point-like quark the virtual correction has the Regge trajectories as coefficients of the rapidity logarithms and contains no $\frac{67}{9}$. The comparison of the finite part of the virtual correction with the one-loop vertex is not completed.
\item Power counting of all contact kernels at all coincidences: all integrable (Section~6.8, Remarks (ii), (iii)); the contact terms do not cancel among the rows (Remark (i)).
\item Prescriptions: the $i\pi\theta$ terms do \emph{not} cancel; we now give both the principal-value (standing-wave) and the out-state ($i\varepsilon$) result, with the difference $\Delta N$ explicitly (Section~6.5, Eqs.~(6.18), (6.19); momentum-space form Eq.~(F.20)). Instantaneous terms: they come from $p^+\sim k^+$, not from the cut-off (Section~6.6).
\item Conditional convergence: a lemma on the uniqueness of the measured-momentum Fourier integrals and an absolutely convergent grouping (Section~6.7, paragraph ``Convergence'', Eq.~(6.24)).
\item Scheme (CDR) stated, with the shift to a dimensional-reduction-like scheme (Section~5.2); BLM-type identification of the $\beta_0$ terms in pure Yang--Mills (Section~5.2, Remark); result versus grouping in Rows~I.II/I.V (Section~5.2, Eq.~(5.7)); separate $\mu_R$, $\mu_F$ (Eq.~(5.20)).
\item Collinear cross-check: standard $e^+e^-\to gX$ and $H\to gX$ coefficient functions reproduced exactly with our convention (Section~5.3).
\item Cubic terms: uniqueness of the split (Section~6.8), exact evaluation on a quark and a dipole (new Section~7.5), cancellation of the four-Wilson-line invariants (Eq.~(6.59)), comparison with Li--Skokov (Section~7.1).
\item All red internal notes, and the macro that produced them, are removed; open items are listed in Section~8.1.
\end{itemize}

\paragraph{On ``What was checked and found consistent''.} We are grateful for these independent checks, in particular of the normalisation of the LO spectrum, of the one-loop wave function and of the pole table. Nothing in the revision changes the items listed there.

% ======================================================================
\section*{Major issues}

% ----------------------------------------------------------------------
\point{1. The manuscript states that the result is unfinished.}
\referee{Sections 6.7, 6.8, 7.6 and 8 list open items: the transition region (C10); integrability of the contact kernels, checked for only two kernels; whether the contact terms cancel between rows; the inclusive check of $\frac{67}{9}-\frac{\pi^2}{3}$; any comparison with the known dilute limit. A contact term at $x=y$ produces $K(x,x,z)=1/(x-z)^2$, which is not integrable at $z\to x$.}

\emph{(a) Contact kernels; the referee's $K(x,x,z)$.}
\done We power-counted every contact line of the colour-singlet equations (the 31 printed lines of the $(2\rho)$ equations of Rows~I.III, I.V, II.I, III.II, III.III, III.V, III.VI, decomposed into 47 kernel$\times$colour units, plus the two lines of Row~I.II, $N_{2b}$) at \emph{every} coincidence of its integration points: 1222 clusters of two to six points, 705 of them containing the merged charge point, and, for a local two-point function, 329 clusters at the triple coincidence of the charges. The exponents were fitted at scales $10^{-2}$ to $10^{-5}$ for three base configurations per line and 16 orientations, with two smooth, non-uniform $SU(3)$ Wilson-line fields (two clusters that lose all digits in double precision were redone with 60 digits).
\result With the colour factors as printed, every line is \emph{absolutely} integrable at every coincidence; the minimal margin is $0.90$ powers of the scale ($\log(1/r)/r$ at a pair). The referee is right that $K(\v x,\v x,\v z)=1/(\v x-\v z)^2$ occurs and is not integrable by itself. It occurs in exactly one kernel, $\mathfrak K^{(z)}_\delta|_{x'\to x}$ of Rows~III.III, III.V and III.VI. In all three lines the colour factor vanishes identically at $\v z=\v x$ by $\mf F(a,a,c)=N_c\mf D(a,c)$ (e.g.\ for III.III: $-N_c\mf D(x,y)-N_c\mf D(z,w)+\mf F(x,y,z)+\mf F(x,z,w)\to0$), so that each line is integrable for any target with H\"older continuous Wilson lines: the term is the commutator remnant of a colour structure that vanishes when the soft gluon sits on the charge. In all other lines the kernel alone is integrable. The large-distance tails of the contact lines cancel pair by pair, as at separate points.
\where Section~6.8, Remarks (ii) and (iii); Section~7.10 item (iv).
\open The power counting is numerical (exponent fits at generic configurations) plus analytic arguments; it is not a proof for every degenerate configuration. The kernel of Row~I.II $N_{2b}$ has no coordinate form; for it we have an analytic estimate only (homogeneous of degree $-1$, $\log r/r$ at worst).

\emph{(b) Do the contact terms cancel between rows?}
\done Lines can cancel at the integrand level only if they share kernel, contact and phase. Such lines exist only within Rows~III.III, III.V, III.VI and between Rows~II.I and III.II; we summed their colour factors.
\result They do \emph{not} cancel: the sums are non-zero (with random $SU(3)$ Wilson lines of modulus $0.4$--$11$, comparable to single lines, $1$--$15$); in II.I$+$III.II the quadrupoles $\mf Q$ cancel but dipoles and $\mf F$ remain. A cancellation against the contact terms of the three-$\rho$ evolution terms is excluded, since those multiply $\log(\vee/k^+)$ or $\log(k^+/\Lambda)$. A cancellation is not needed, by (a). In the dilute limit the commutator (contact) parts of the real rows are in fact \emph{required} to reproduce the known two-gluon production cross section (Major~2).
\where Section~6.8, Remark (i).

\emph{(c) Transition region of (C10).}
\done The three rows of (C10) are, at fixed $t=p^+/k^+$, $|\mathcal A_B-\mathcal A_{CA}|^2$ with the measured gluon Fourier transformed. The large-$|\v z|$ expansion of the Fourier-transformed amplitude is determined by the singular points of the outer-region function (measured gluon at the projectile, at the other gluon, or in the collinear configuration); the outer region, in which only part of the pair is far, is smooth and contributes only at higher order (integration by parts). We verified this numerically on the full $\v w$ plane (one charge, colour stripped, target uniform at large distance).
\result The residual falls like $\log Z/Z$ for $t=1$ ($0.075\to0.018$ from $Z=10$ to $80$, in units of $1/|\v k|$) and like $1/Z$ for $t=0.3$ ($Z\times$residual $=4.1,4.6,4.7$); a contribution of the transition region at the order of the tail would be a $Z$-independent residual of size $|c_{\rm joint}|=3.9$ resp.\ $7.8$. The transition region does not contribute.
\where Section~6.7, condition (C10); Sections~5.5 and 6.9.
\open For a colour-neutral projectile and a target that is not uniform at large distance the statement rests on the argument given in (C10), not on a numerical check.

\emph{(d) Inclusive check of $\frac{67}{9}-\frac{\pi^2}{3}$.} See the error item on Section~7.6 below: the claim is withdrawn, and the constant of Row~I.II is shown to be a property of the grouping of Rows~I.II and I.V.

\emph{(e) Dilute limit.} See Major~2.

\emph{Finiteness.} With (a)--(c) and the answers to Majors~5 and~6, the finite terms are shown to be finite and integrable in the grouping given; none of these checks required a correction of the finite terms. The manuscript no longer describes the result as unfinished; what is not done is listed in Section~8.1.

% ----------------------------------------------------------------------
\point{2. No comparison with any independent NLO result.}
\referee{The dilute--dilute limit must reproduce the one-loop central-emission (Lipatov) vertex plus real $gg$ production in $k_T$ factorisation. This is the decisive test.}
\done We took the dilute--dilute limit ($U=1+u$, $O(u)$ in the amplitude and in the conjugate amplitude, projectile two-point function, contact terms from $[\rho,\rho]$ kept) at the level of the amplitudes and coefficients of the paper, separately for the real emission (Groups~II and III) and for the virtual correction (Group~I).
\result \emph{Real emission.} At $O(u)$ the anticommutator parts of the two-gluon amplitude cancel exactly, the amplitude becomes linear in $\rho$, and the real part of the NLO result is \emph{exactly} the Fadin--Lipatov quasi-multi-Regge two-gluon production cross section, in the form $A_{\rm gluons}$ of Bartels, Sabio Vera and Schwennsen, Eq.~(7.2): the ratio is $1$ to $2\times10^{-15}$ at random kinematics, for all $0<x<1$, for $SU(2)$ and $SU(3)$, and for $d_\perp=1.7,3,6$ transverse polarisations against their $(D-2)$ terms; there is no free normalisation. Without the commutator (contact) parts, or without $\mathcal A_{CA}$, the ratio is $0.2$--$10$. Both strong-ordering limits give the real part of the LO BFKL kernel with $\bar\alpha_s=\alpha_sN_c/\pi$, on $\phi_P$ and on $\phi_T$. Since $dp^+/p^+=dx/(x(1-x))$ is the rapidity measure of the second gluon at fixed $k^+$, our finite real remainder coincides identically with BSV's QMRK-minus-MRK remainder for $s_\Lambda\to\infty$ (at three random points $-2.62950455$, $-0.66773731$, $-6.13068770$ against $-2.62950704$, $-0.66773773$, $-6.13068408$ at $s_\Lambda=10^6$, differences falling like $1/s_\Lambda$): for the real emission no scheme conversion is needed. The coefficient of $\log\mu^2$ and the scale of the coupling agree with FL and BSV.

\emph{Virtual correction, point-like quark.} For $\rho^a=t^a\delta^{(2)}(\v x)$ we evaluated all of Group~I in momentum space ($d_\perp=2-\epsilon$, $\MS$, $\mu=|\v k|$, two kinematic points). The result, Eq.~(7.3), is
\begin{align*}
\frac{\text{Group I}}{\rm LO}&=2\omega(\v q_a)\log\frac{\vee}{k^+}+2\omega(\v q_b)\log\frac{k^+}{\Lambda}+\frac{\alpha_sN_c}{\pi}\Big[\frac2\epsilon\log\frac{k^+}{\Lambda}+\log^2\frac{k^+}{\Lambda}\Big]+C(\v k,\v l)\,,\\
2\omega(\v q)&=\frac{\alpha_sN_c}{\pi}\Big[\frac2\epsilon-\log\frac{\v q^2}{\mu^2}\Big].
\end{align*}
The coefficient of $\log(\vee/k^+)$ is the trajectory $2\omega(\v q_a)$ exactly ($-0.0592216$ against $-0.0592216$; $0.0677534$ against $0.0677533$); the coefficient of $\log(k^+/\Lambda)$ is $2\omega(\v q_b)$ to $2\times10^{-4}$ ($0.04145$ against $0.04139$ at a point where $2\omega(\v q_a)=-0.059$), but only in the sum of the rows. In addition there is the $\v l$-independent soft--collinear term $\frac{\alpha_sN_c}{\pi}[\frac2\epsilon L+L^2]$. The $\frac{11}{3}$ poles cancel between Rows~I.II and I.V, and the unresolved part of Row~I.V cancels the whole self-energy part of Row~I.II, $\xi$ by $\xi$ to $10^{-13}$, including the $\log(\xi(1-\xi))$ that produces $\frac{67}{9}$: Group~I contains no $\frac{67}{9}$. The finite remainder is a function, $C=0.055$ and $0.305$ times LO at the two points. The principal-value and the out-state prescriptions give the same result in this limit.
\where New Section~7.2, Eqs.~(7.1)--(7.3); Section~7.10 item (v); Section~5.2 (identity (5.7)).
\open The decisive part of the referee's test, the comparison of the finite part of the virtual correction with the one-loop vertex, is \emph{not} completed. FL/BSV regulate the soft--collinear region by $\epsilon^{-2}(\boldsymbol\Delta^2/\mu^2)^\epsilon$, we by the plus-momentum cut-off; the two virtual parts therefore differ already in their pole terms by a function of $\boldsymbol\Delta$ that is not of the form $F_P(\v q_a)+F_T(\v q_b)$, and a scheme-independent comparison needs virtual plus real with matched regulators. In particular the real emission with the cut $p^+>\Lambda$ is not uniform at $|\v p|\lesssim(\Lambda/k^+)|\v k|$; whether this region cancels the extra term $\frac{\alpha_sN_c}{\pi}[\frac2\epsilon L+L^2]$ has not been computed. Also open: the conversion of the virtual trajectories between plus-momentum cut-offs and BSV's energy scales, Group~I for a general projectile, the dilute limit of the coordinate-space finite terms as printed. No disagreement with FL/BSV was found; we do not claim agreement for the finite part of the virtual correction (Section~8.1 (i)).

% ----------------------------------------------------------------------
\point{3. The running-coupling conclusion is not established.}

\referee{In pure Yang--Mills, $\beta_0/(2N_c)=11/6$ is numerically identical to integrals of the splitting function \dots\ ``$\beta_0$ terms'' cannot be separated from these without an $N_f$ (quark-loop, BLM-type) calculation.}
\done The calculation is, by design, for pure Yang--Mills theory, and the result remains $N_f=0$; we do not compute the $N_f$ rows of the result. We did, however, carry out the BLM-type identification the referee asks for, at the level of the kernels: with the vertices in the paper's normalisation (the triple-gluon matrix element is exactly minus the covariant vertex contracted with light-cone polarisation vectors, the quark vertex in the same normalisation), the angular-averaged squared vertices give, for every $\xi$, $\frac{n_fT_R\sum|V_{q\bar q}|^2}{\frac12N_c\sum|V_{gg}|^2}=\frac{n_fT_R}{N_c}\frac{\xi^2+(1-\xi)^2}{K_{gg}(\xi)}$, i.e.\ a quark loop adds $\frac{n_fT_R}{N_c}[\xi^2+(1-\xi)^2]$ to the self-energy kernel $K_{gg}=\xi(1-\xi)-2+\frac{1}{\xi(1-\xi)}$, and the corresponding $q\bar q$ pair crosses the shock wave in Row~I.V.
\result The $\frac{11}{3}$ of Row~I.II becomes $\beta_0/N_c$; the $-\frac{11}{6}$ of $N_{2b\,\rm UV}$ of Row~I.V becomes $-\beta_0/(2N_c)$; the $\frac{11}{6}$ of the finite kernel $\mathfrak K^{\rm V}_{2b}$ becomes $\beta_0/(2N_c)$; the $\frac{11}{12}$ of the normalisation, of the counterterm and of the RC term become $\beta_0/(4N_c)$; the $\frac{\beta_0}{2}\delta(1-\xi)$ becomes the full $\frac{\beta_0}{2}$. In contrast, the $-\frac{11}{6}$ tail coefficient of Eq.~(5.26) (old (5.24)) is the double-plus integral of the real $\mathcal P_{gg+}$ and stays $-\frac{11}{6}$ for any $N_f$. The Row~I.II constant would become $\frac{67}{9}-\frac{\pi^2}{3}-\frac{20}{9}\frac{n_fT_R}{N_c}=\frac{2}{N_c}K_{\rm CMW}$. The pole cancellations and the pole table hold for any $\beta_0$.
\where Section~5.2, ``Remark (pure Yang--Mills; identification of the $\beta_0$ terms)''; Section~7.9.
\open The $N_f$ terms of the result are not computed (not part of this pure Yang--Mills calculation).

\referee{The claimed exact cancellation of the $B(x-w)$ logarithms between Rows I.II and I.V assumes a point-like pair \dots\ ``The emission coupling does not run with $x-w$'' is therefore a choice of grouping, not a result.}
\done We agree, and we have made the distinction explicit. We found an exact LCPT identity: with $a$ the energy of the emitted gluon, $D$ the relative energy of the pair and $b=a+D$, $\frac{1}{ab}+\frac{1}{bD}=\frac{1}{aD}$, whose transverse integral is scaleless (Eq.~(5.7)). Hence, with the colour of the pair evaluated at its centre of mass, the self-energy part of Row~I.II (pole, $\mathcal B$ logarithm, soft end points and the constant $\frac{67}{9}-\frac{\pi^2}{3}$) is cancelled identically by the unresolved part of Row~I.V; this was confirmed independently in momentum space in the dilute limit (to $10^{-13}$).
\result The cancellation of the poles and of $\log\mu^2$ is a result. That no term $\frac{\alpha_s\beta_0}{4\pi}[\mathcal B(\v x-\v w)+\mathcal B(\v x'-\v w')]\times$LO remains explicitly is a statement about the grouping (colour of the pair at one of its gluons). The resolving region sits in the finite terms of Row~I.V with the coefficient $2L-\frac{11}{6}$; for an emitter distance larger than $1/Q_s$ it contains, by power counting, a logarithm of $|\v x-\v w|Q_s$ with this coefficient. The coupling of the emission vertex therefore does not run with $\v x-\v w$ only as long as the pair does not resolve the target within this distance.
\where Section~5.2, paragraph ``Result and grouping'' (replacing the old sentence); Section~8.
\open The resolving window was not evaluated with a target model (Section~8.1 (v)).

\referee{Horowitz--Kovchegov are described as running ``with the collinear resolution scale''. As recalled, their scale $\Lambda_{\rm coll}$ depends on $q$ and $k-q$; this description should be checked against that paper.}
\done We re-read the paper (arXiv:1009.0545; equation numbers of the arXiv version).
\result The recollection in the report is, we believe, not correct. In HK, Eq.~(3.10), $\Lambda_{\rm coll}$ is a fixed upper cut-off on the virtuality $k^2$ of the produced-gluon line, i.e.\ on the invariant mass of the collinear pair, described as the minimum momentum scale resolved by the detector; it does not depend on $\v q$ or $\v k-\v q$. The produced gluon carries $\alpha_s(\Lambda^2_{\rm coll}e^{-5/3})$ (Eqs.~(3.11)--(3.16)). The dependence on $\v q$ and $\v k-\v q$ is in the other couplings, $\alpha_s(\v q^2e^{-5/3})$, $\alpha_s((\v k-\v q)^2e^{-5/3})$, and in the complex scale $Q(\v q,\v k-\v q)$ of Eqs.~(3.31), (3.32). Our wording ``runs with'' was nevertheless imprecise, and we now say that the coupling is evaluated at a fixed cut-off on the invariant mass of the collinear pair, which plays the role of the factorisation scale.
\where Sections~5.4 and 7.9 (fifth bullet).

\referee{Showing that the residual $\mu$-dependence comes only from the $\delta(1-\xi)$ term of $P_{gg}$ is an add-and-subtract that amounts to setting $\mu_R=\mu_F$, which the manuscript never separates.}
\done We now separate the scales. With $g_0^2=\mu_R^\epsilon g^2(\mu_R)[1-\frac{\alpha_s\beta_0}{4\pi}\frac2\epsilon]$ and $D_0=[\delta(1-\xi)+\frac{\alpha_s}{2\pi}\frac2\epsilon\mathcal P_{gg}]\otimes D(\mu_F)$ the two halves of the scaleless normalisation are cancelled separately.
\result LO$+$RC carries $g^2(\mu_R)$ and $\log((\v y'-\v y)^2\mu_R^2/c_0^2)$, the DGLAP term (with its $\delta(1-\xi)$ term) is at $\mu_F$, and their $\beta_0$ part is $+\frac{\alpha_s\beta_0}{4\pi}\log\frac{\mu_R^2}{\mu_F^2}\times$LO; the result is independent of $\mu_R$ and $\mu_F$ to this order. The referee is right that our earlier statement holds only for $\mu_R=\mu_F$; it is now restricted accordingly.
\where Section~5.4, paragraph ``Separate renormalisation and factorisation scales'', Eq.~(5.20); Section~6.1.

% ----------------------------------------------------------------------
\point{4. The scheme of the finite constants is not specified.}
\referee{State whether gluon polarisation indices are continued to $d_\perp=2-\epsilon$ (CDR) or kept two-dimensional \dots\ otherwise the constant shifts by a term of order $N_c/6$.}
\done We traced the polarisation sums through all rows with poles and state the scheme.
\result The calculation is in conventional dimensional regularisation: polarisations are continued to $d_\perp=2-\epsilon$ wherever a polarisation sum multiplies a pole (Rows~I.II, I.V, II.IV); the finite terms are locally integrable and taken at $d_\perp=2$; the coupling is $\MS$ and the fragmentation function the $\MS$ number density, both in this scheme; $\frac{67}{9}-\frac{\pi^2}{3}$ and Eq.~(2.16) are CDR values. With two-dimensional polarisations the Row~I.II constant would be $\frac{64}{9}-\frac{\pi^2}{3}$, a shift $-\frac{\alpha_sN_c}{6\pi}\times$LO, which is the shift of order $N_c/6$ the referee mentions; it is cancelled by the opposite shift of Row~I.V (as it must be by the identity (5.7)); Row~II.IV shifts by $-\mathcal C\int d\xi\,2N_c\xi(1-\xi)\Pi(\xi)$; and with $\alpha_s^{\rm DR}=\alpha_s^{\MS}(1+\frac{\alpha_sN_c}{12\pi})$ and the corresponding conversion of the fragmentation function (whose $\delta(1-\xi)$ coefficient $\frac{N_c}{6}$ is fixed independently by the momentum sum rule) the result is unchanged. Independently, in the dilute limit our two-gluon amplitude has the CDR $(D-2)$ structure of FL/BSV for any $d_\perp$.
\where Section~5.2, paragraph ``Scheme''; Section~7.2.
\open The sign of the Row~II.IV shift was obtained from the general $d_\perp/d$ rule and the closure of the conversions, not by re-deriving the row (stated in the text).

\referee{The claim that Chirilli--Xiao--Yuan and Altinoluk et al.\ use a convention that violates the momentum sum rule is strong. The remedy is the check the red note on p.~31 says was ``not done'': reproduce a standard NLO $pp\to gX$ (or $e^+e^-\to gX$) coefficient with the same conventions.}
\done We computed two standard one-hadron-inclusive gluon coefficient functions with exactly our conventions (CDR, $\MS$ coupling, pure-pole $\MS$ number-density fragmentation function, counterterm with the $d$-dimensional Jacobian $\xi^{-d}$), in exact arithmetic: $e^+e^-\to gX$ (Altarelli et al., Furmanski--Petronzio, Nason--Webber, as given by Kretzer) and, as the direct analogue of the $g\to g$ final-state splitting in pure Yang--Mills, $H\to gX$ (Almasy--Moch--Vogt).
\result Both are reproduced \emph{exactly} (pole and finite part). With the two-dimensional Jacobian $\xi^{-2}$, which is what produces the CXY form, both come out wrong by exactly $-\frac{\alpha_s}{2\pi}P(z)\log z^2\times$Born, the $z$-structure (and sign) of our $\Delta$, Eq.~(5.13). Our convention is that of Catani--Seymour (their Eq.~(6.17), $dz/z^{2-2\epsilon}$). We confirmed from the $d$-dimensional formulas of CXY and AABKL (App.~C) that their counterterm corresponds to $\xi^{-2}$. We have reworded the claim: CXY and AABKL use standard (fitted) $\MS$ fragmentation functions, which do obey the momentum sum rule; the precise statement is that their collinear term is the coefficient for the functions $D^{\rm CXY}=[\delta+\frac{\alpha_s}{2\pi}\mathcal P_{gg}\log\xi^2]\otimes D$, which do not, since $\int_0^1d\xi\,\xi\,\mathcal P_{gg}\log\xi^2=4N_c(\frac{65}{144}-\zeta_2)=-14.3225$ ($N_c=3$); used with standard fragmentation functions it misses $\Delta$ at $O(\alpha_s)$.
\where Section~5.3 (paragraphs on the definition of the fragmentation function and ``Remark: the CXY convention''; the red note is removed); Section~7.9 (third bullet).
\open A comparison with the NLO hadroproduction coefficient $pp\to gX$ and a collinear limit of the hybrid formula were not done; the bare final-state collinear terms of CXY and ours are the same $d$-dimensional integral, so such a comparison would re-test only the counterterm. The two checks are angle-integrated; for the transverse-momentum Jacobian they rely on universality and $d$-dimensional collinear kinematics (Catani--Seymour's derivation). Section~8.1 (vi).

% ----------------------------------------------------------------------
\point{5. Prescription dependence.}
\referee{Vanishing energy denominators in $B_3$ are taken as principal values. The physical (adiabatic, $i\epsilon$) prescription for the in-state differs \dots\ Show that the $i\pi\theta$ terms also cancel in the combined kernel (6.19), not only the pole.}
\done We derived the $i\varepsilon$ form from the adiabatic (Gell-Mann--Low) construction. The only vanishing energy denominators are $E_g(k)-E_g(p)$ of $\mf B_3$, which enters only the final-state dressing and hence only Row~I.III (the in-state has no vanishing denominator at this order; the relevant state is the out-state of the observed gluon). We evaluated the residues of the separate terms on the energy shell and summed them.
\result The referee's concern is justified: the $i\pi\theta$ terms do \emph{not} cancel. In the combined kernel they combine, exactly, into a constant: $\mathcal M^{ij}|_{i\eta0}=\mathcal M^{ij}|_{\rm PV}-i\eta\pi\,\mathfrak B^{ij}$, where $\mathfrak B^{ij}=[\cos2\phi\,\delta^{ij}+\sin2\phi\,\epsilon^{ij}]/[2(\v z-\v w)^2]$ is literally the brace that multiplies $\log\frac{(x-z)^2}{(x-w)^2}$, i.e.\ $\log\frac{(x-z)^2}{(x-w)^2}\to\log\frac{(x-z)^2}{(x-w)^2}-i\eta\pi$ ($\eta=+1$ out-state, $-1$ in-state dressing of the final state, $0$ principal value). The pole on the line cancels for every $\eta$. The on-shell part of $\mf B_3$ is anti-Hermitian (checked on the shell to $4\times10^{-13}$), so the difference conserves the total number of gluons, but it changes the spectrum, and the complex conjugate does not cancel it. The difference is the finite term $\Delta N$, Eq.~(6.19): a forward Coulomb logarithm $\log(k^+/\Lambda)$ at $\v z\to\v w$, multiplying $\mathcal A^\dagger_{\rm LO}f\cdot\bar Q\,\mathcal A_{\rm LO}$, cancels exactly against the complex conjugate, and the remainder is integrable. Its symmetrised three-$\rho$ part is even and its commutator part odd in $\v k$ (the opposite of the principal-value terms). In momentum space $\Delta N$ is the on-shell Coulomb rescattering of the LO gluon on the rotated charge, with the on-shell vertex $n_{ji}=\frac{|\v k||\v p|}{(\v k-\v p)^2}R(2\theta)_{ji}$ (Eqs.~(F.19), (F.20)). In the dilute point-quark limit $\Delta N=0$. Following this, we now present both prescriptions: the finite terms as printed are those of the principal-value (standing-wave) dressing, and those of the out-state dressing are obtained by adding $\Delta N$ with $\eta=1$.
\where Section~6.5 (rewritten: ``Principal value'', ``Out-state dressing'', ``Which prescription''), Eqs.~(6.18), (6.19); condition (C6); Appendix~F.4; Sections~2.3, 6.1, 6.2, 6.9, 7.10 and Appendix~E (cross-references).
\open Which prescription is the physical one within the eikonal approximation, where the static charges interact with the produced gluon for an infinite light-cone time, cannot be decided within this calculation; we give both. The size of $\Delta N$ for a realistic configuration has not been evaluated, and the coordinate and momentum-space forms were matched only through the forward coefficient and the tensor structure (Section~8.1 (iii)).

\referee{The instantaneous Group III terms converge only as distributions, $\pi\delta(s)-i\mathcal P(1/s)$, as $\vee\to\infty$. They come from $p^+$ near $\vee$, where the eikonal approximation fails, so their finite parts are not under control.}
\done We analysed the $p^+$ distribution after the transverse ($\v y'$) integral, and the dependence on the shape of the soft/valence separation.
\result Pointwise in the transverse coordinates the integrand indeed does not decay in $p^+$, but after the $\v y'$ integral it does: the distribution falls like $(k^+/p^+)^2$, because the bracket vanishes at coinciding points (for a H\"older-$\alpha$ target like $(k^+/p^+)^{1+\alpha}$). Hence $p^+>P$ contributes $O(k^+/P)$ [$O((k^+/P)^\alpha)$], and the bulk comes from $p^+\lesssim k^+\max\{1,1/(|\v k|\ell)\}$, $\ell$ the shortest transverse scale ($1/Q_s$); the eikonal approximation requires $\vee$ to be large compared to this, not more. With any smooth separation $f(p^+/\vee)$ the finite parts have the same limit (deviations $O(k^+/\vee)$); the only dependence of the whole result on $f$ is the universal shift $\log(\vee/k^+)\to\log(\vee/k^+)+c_f$ of the BFKL logarithm, $c_f=\int_0^\infty\frac{du}{u}[f(u)-\theta(1-u)]$ ($0,-\gamma_E,-\gamma_E/2,0$ for four shapes), i.e.\ a change of the rapidity factorisation scale. At finite $\vee$ all integrals converge absolutely, and the distributional form is the limit. On model brackets with the right short-distance behaviour, 33--89\% of the limit comes from $p^+<10k^+$ and 78--99\% from $p^+<100k^+$.
\where Section~6.6, new paragraph ``Where in $p^+$ the instantaneous terms come from''; condition (C7).
\open The checks use model brackets, not the actual one with Wilson lines; that the same $c_f$ appears in every row is argued from the $1/p^+$ form of the integrands, not checked row by row (Section~8.1 (vii)).

% ----------------------------------------------------------------------
\point{6. The finite terms are only conditionally convergent.}
\referee{Their values depend on the order of integration and on angular prescriptions: conditions (C3), (C7), (F.12), and the two integrals in (6.15) \dots\ Present such a grouping, or prove that the prescriptions cannot shift the result by a finite amount.}
\done We prove a uniqueness lemma for the integrals that carry the measured phase and give an absolutely convergent grouping; each of the four items was checked with several regulators and orders.
\result \emph{Lemma} (Eq.~(6.24)): for a locally integrable, tempered $f$ whose Fourier transform is continuous at the measured momentum, $\lim_{\epsilon\to0}\int d^2r\,e^{-i\v k\cdot\v r}\phi(\epsilon\v r)f(\v r)=\hat f(\v k)$ for \emph{every} regulator $\phi$ with $\phi(0)=1$ and integrable Fourier transform (exponential or Gaussian of any direction dependence, smooth cut-offs of any shape): the left-hand side is a convolution with an approximate identity. ``Angular integral first'' and ``radial integral first'' are special cases. With these integrals done first, all remaining integrals converge absolutely in the grouping of Section~6.7.
(C3) is the definition of the measured momentum, not a choice: at fixed $\v w$ the angular mean of the $1/\v z^2$ tail is $\Kv^i(\v x-\v w)[\gamma_E+\log(|\v k||\v x-\v w|)-\log2]/\v z^2\neq0$, so the other order does not exist; there is no second finite value.
The two integrals of old Eq.~(6.15) (now (6.16)) are Fourier integrals of the type of the lemma and each is unique; what is conditional is only the subsequent $\v z$ integral of each term alone, and in the grouping they cancel for every $\theta$. Numerically, with the isotropic factor $e^{-\epsilon r}$ and two anisotropic ones (non-circular cut-off regions), the deviations are $\le1.2\times10^{-3}$ at $\epsilon=10^{-4}$ and fall like $\epsilon$.
(F.12): at fixed $a=\v k\cdot(\v y-\v y')$ the $\xi$ integral is $(1+ia)E_1(ia)-e^{-ia}$, bounded by $C(1+|\log|a||)$ uniformly in a lower cut-off, so the two orders agree (on a model with the structure of the term: $-3.0246763985$ against $-3.0246763984$).
(C7): the only non-absolutely convergent piece is independent of $s$ and is annihilated by $\int d\kappa\,e^{i\kappa}[\pi\delta(\kappa)+i\mathcal P\frac1\kappa]=0$; any large-distance cut-off, symmetric or not, gives the same result (three orders, including an asymmetric cut $(-U,3U)$, agree to $3\times10^{-4}$, the quadrature accuracy).
In doing this we found one genuinely conditionally convergent \emph{local} point not on the referee's list: the collinear point $\v z\to\v x$ of $\mathfrak K^{\rm V}_{2b}$ (Row~I.V), a pure second harmonic $a_2\cos2(\varphi-\varphi_0)/|\v z-\v x|^2$ (mean and first harmonic vanish). It is zero for any rotation-invariant cut-off, which is the regularisation in which its $\frac{11}{6}$ subtraction was obtained; we give an absolutely convergent subtracted form.
\where Section~6.7, conditions (C3), (C7) and paragraph ``Convergence''; Eq.~(6.16) and the text after it; Appendix~F after Eqs.~(F.5) and (F.12).
\open At the Row~I.V collinear point, continuing the transverse indices to $d_\perp=2-\epsilon$ gives the traceless tensor a trace $O(\epsilon)$ that multiplies the $1/\epsilon$ of the collinear subtraction and could leave a finite local term; whether it is present is not decided (Section~8.1 (vi)).

% ----------------------------------------------------------------------
\point{7. The physics claims about the cubic terms need support.}
\referee{The split of the cubic terms into an odd-in-$k$ part $\propto o$ and even ``contact terms'' depends on operator ordering.}
\done We examined where the ordering comes from and what the split means independently of it.
\result The ordering is not a choice: it is the order in which the coefficients of $\bar\Omega^\dagger\Omega$ and of the conjugate amplitude act, fixed by the light-cone time ordering (Eq.~(2.10); e.g.\ $[\mf A^{(3)}]_{\rho\rho\rho}=\mf A\,\mf N_2$, verified against an exact construction of $\Omega$). The decomposition is an operator identity, and its symmetrised part is the \emph{same operator for every ordering}; only the contact terms depend on the ordering (on a dipole the six orderings of one structure give the same symmetrised part to $3\times10^{-17}$). Under (A1), (A2) the two parts have an ordering-free meaning: since charge conjugation reverses $o$ and leaves $n$ unchanged, the symmetrised part is $\frac12[\sigma(P)-\sigma(P^{\rm C})]$ and the contact part $\frac12[\sigma(P)+\sigma(P^{\rm C})]$, and by C-invariance of the result (verified exactly for all thirteen structures) these are also the parts odd and even under $V\to V^*$. The split is therefore unique.
\where Section~6.8, paragraph ``The split is unique''.

\referee{Its physical content should be shown on a case with a unique answer --- a single quark or dipole computed exactly in charge space --- and compared with Li--Skokov and McLerran--Skokov.}
\done We evaluated all thirteen cubic colour structures exactly in charge space (explicit $SU(2)$, $SU(3)$, $SU(4)$ matrices) for a quark, an antiquark, a dipole and the conjugate dipole, all assignments of the charge points, random Wilson lines; and we compared with Li--Skokov (arXiv:2104.01879), McLerran--Skokov and Kovchegov--Skokov.
\result \emph{Quark.} $\langle\rho^a\rho^b\rho^c\rangle=\frac{1}{4N_c}(d^{abc}+if^{abc})$; the symmetrised part is the $d^{abc}$ half, $o=C_F\frac{N_c^2-4}{2N_c}$ ($\frac{10}{9}$ for $SU(3)$), the contact part the $if^{abc}$ half; for the antiquark the first changes sign, the second does not. \emph{Dipole.} $o\propto s_is_js_k$, $n\propto s_is_j$ ($s=\pm1$ at quark/antiquark): the $\v k$-odd cubic term is the part of the spectrum that distinguishes the quark from the antiquark end and needs a C-odd target; the contact part survives the orientation average. Exact values agree with the split to $5\times10^{-14}$. Neither part is enhanced for a quark or a dipole (colour weights relative to LO for a dense $SU(3)$ target: about $0.12$ symmetrised, $0.8$ contact); ``first saturation correction'' applies to the symmetrised part only for a coherent dense projectile with large $o$.
\emph{Li--Skokov, McLerran--Skokov.} The colour structures agree: the two-Wilson-line term of LS has zero $d^{abc}$ projection, as ours; the odd part needs three Wilson lines and a C-odd target; and we found that our four-line invariants $\mf Q_d$ cancel between Rows~II.I and III.II (new, Eq.~(6.59), to $1.5\times10^{-15}$), so that our odd part has at most three Wilson lines, as in LS and MS. The parities in $\v k$ do \emph{not} agree. LS have a $\v k$-even part configuration by configuration, which our symmetrised cubic terms (principal value) do not have. We tested whether the out-state term $\Delta N$ (Major~5), whose symmetrised part is even, accounts for it: it does not (kernel-level residual above $90\%$; on a lattice with dense random $SU(3)$ Wilson lines the ratio of $\Delta N$ to LS's even term varies between $-0.25$ and $3.9$); nor does the principal-value result reproduce only their odd part (in the two-line sector LS's term is even, ours odd, with different kernels). We traced the difference to a relative phase $i$: the transverse kernel of the classical $O(g^3)$ Weizs\"acker--Williams field of LS appears in our wave function (Eq.~(B.16), to $10^{-12}$), but with a factor $i$ relative to the LO coefficient, and the same holds for the two-charge structures of $\mf A^{(3)}$, the $f\rho$ part of $\mf B_2$ and the rescattering in $\mf B_3$, relative to LS and to Kovchegov--Skokov's LCPT with retarded propagators; for the gluon--gluon merging after the collision ($\mf C_1$) the assignments agree.
\where New Section~7.5; Section~7.1, paragraphs ``Odd terms'' and ``Comparison with Li and Skokov''; Section~6.8 (Eq.~(6.59)).
\open The parity difference with Li--Skokov is not resolved. A possible origin is the prescription for $1/k^+$ (retarded versus symmetric $x^-$ profile of the field), which we have not checked; LS's even part has not been confirmed independently (LS describe it as computed for the first time); a kernel-level comparison of the odd three-line terms needs all cubic rows in momentum space and was not done. We state this as an open point (Section~8.1 (ii)) rather than claim agreement.

% ======================================================================
\section*{Errors and inconsistencies}

\point{p.~64, Section 7.3: the dipole emission probability is written as $(g^2/4\pi^4)C_FK_D$; it must be $g^2C_F/(4\pi^3)$.}
\result The manuscript does not contain $g^2/(4\pi^4)$ there. We checked the source of all versions and the submitted PDF: the text reads ``the probability to emit the additional gluon at $\v z$ is $\frac{g^2}{4\pi^3}\log\frac{\vee}{k^+}C_F\,K_D$'', i.e.\ $\frac{\alpha_sC_F}{\pi^2}K_D$, which is what the referee asks for. The likely source of the confusion is the undefined factor ``(LO prefactor)'' in old Eq.~(7.2), which is $\frac{\alpha_s}{4\pi^4k^+}$. We re-derived the probability from our Eq.~(2.14) ($-2\mf N_2$ for the dipole) and verified the dipole equation exactly at finite $N_c$: with $\frac{g^2C_F}{4\pi^3}$ the ratio is exactly $1$; with $\frac{g^2}{4\pi^4}$ it would be $1/\pi$.
\where Section~7.4: the probability is now also written as $\frac{\alpha_sC_F}{\pi^2}$, the LO spectrum of the dipole $\frac{\alpha_s}{4\pi^4k^+}\mathcal D_{12}$ is stated, and ``(LO prefactor)'' is replaced by $\frac{\alpha_s}{4\pi^4k^+}$ in Eq.~(7.5) (old (7.2)).

\point{The 11/6 term: Section 5.2 says the 11/6 part of Row I.V cancels Row I.II completely, yet (6.29)/(F.5) still contain an $N_c\cdot\frac{11}{6}$ finite kernel.}
\result There is no contradiction, but the text did not explain it. Row~I.V is split as $N_{2b}=N_{2b\,\rm UV}+N_{2b\,\rm rest}$ by adding and subtracting the small-pair (UV) integrand. The $-\frac{11}{6}$ of $N_{2b\,\rm UV}$ is the dimensionally regularised $\v z$ integral of the subtraction term, and it is this integrated term that cancels Row~I.II. The same term, unintegrated and with the opposite sign, remains in $N_{2b\,\rm rest}$, where it is the $\frac{11}{6}$ of $\mathfrak K^{\rm V}_{2b}$ and makes the collinear point $\v z\to\v x$ integrable; its $\v z$ integral alone diverges. ``Cancels completely'' is true of the integrated subtraction, not of the kernel.
\where Section~5.2 after Eq.~(5.6); Appendix~F after Eq.~(F.5); kernel in Eq.~(6.33) (old (6.29)).

\point{Section 5.2: a leftover pole $\propto\ell_\vee/2$ is removed only by taking $\vee\to\infty$ before $\epsilon\to0$. State this as part of the eikonal approximation rather than as a cancellation.}
\result Agreed, and done. The coefficient $\ell_\vee=\log(1-k^+/\vee)=-k^+/\vee+\dots$ is a power correction; it is the part of the vertex correction (and of the real splitting) in which the parent of the observed gluon is faster than $\vee$ and which is not enhanced at small momentum fraction. Its coefficient is not defined by the effective theory: for a smooth soft/valence separation $f$ it becomes $-\kappa_fk^+/\vee$, $\kappa_f=\int_0^\infty dv\,[1-f(v)]/v^2$ ($1,\sqrt\pi,\pi/2$ for three shapes, checked numerically). We now define the eikonal approximation as the leading term of the expansion in $k^+/\vee$ and $\Lambda/k^+$, with the logarithms kept, at the level of the longitudinal integrands at fixed transverse momenta and fixed $\epsilon$; since the pole factor is a product of a longitudinal and a transverse factor, this is uniform in the transverse loop momentum and, at fixed $\epsilon$, equivalent to $\vee\to\infty$ first. We no longer call it a cancellation.
\where Section~5.2, paragraph ``The balance''; condition (C1); Section~7.9 (last bullet); Section~8.

\point{Section 7.6: the argument that the other half of the constant ``must reside in the collinear finite terms'' is not verified.}
\result We withdraw the argument. Two new results show that it was not only unverified but misdirected. (1) By the identity (5.7), the self-energy part of Row~I.II, which carries the whole constant $\frac{67}{9}$ together with the $\beta_0$ logarithm, is cancelled identically by the part of Row~I.V in which the target rotates the pair at its centre of mass: the constant (7.6) is a property of the grouping of Rows~I.II and I.V, not a constant of the result. This is confirmed in momentum space in the dilute limit with a point-like projectile, where Group~I contains no $\frac{67}{9}$ (Major~2). (2) We evaluated the momentum-weighted collinear terms of Row~II.IV exactly: they give $\frac{\alpha_sN_c}{4\pi}(\frac{65}{18}-\frac{4\pi^2}{3})=-9.548\,\frac{\alpha_sN_c}{4\pi}$ times LO in our scheme (zero in the CXY scheme), not $-4.155\,\frac{\alpha_sN_c}{4\pi}$, plus a linearly polarised term from the tensor term that is not proportional to LO. The standard $H\to gX$ coefficient function, in which $\frac{67}{9}$ sits only in the $\delta(1-z)$ term and the energy-weighted integral is the process-dependent total rate, shows that no general equality should be expected. In the dilute limit the real rows are exactly the QMRK cross section, whose pair-inclusive integral gives FL's $\frac{\alpha_sN_c}{4\pi}(\frac{67}{9}-\frac{\pi^2}{3})$.
\where Section~7.8 (rewritten; the red note is removed); Sections~5.2, 6.8 (text after Eq.~(6.34)), Appendix~F after Eq.~(F.2).
\open The LO-proportional constant left in the finite terms of Row~I.V in our grouping has not been evaluated for a general projectile and target, and the collinear regions of Rows~II.I, III.III and I.V have not been integrated (Section~8.1 (v)).

\point{Section 4.3: the statement that one-step evolution ``cannot produce the C-odd $o$'' \dots\ follows from the C-parity of the Hamiltonian; say so.}
\result Rewritten with C-parity as the reason. We give the transformation ($\rho^at^a\to-(\rho^at^a)^T$, $V\to V^*$, $A^a\to-\eta_aA^a$), show that it preserves all couplings, list the C-parities ($n$, $\mf D$, $\mf F$, $\mf Q$ even; $o$, $\mf G$, $\mf Q_d$ odd), and derive from the factorised form of the logarithms that neither logarithm can multiply $o$. We add one precision: C alone allows terms $o\times$(C-odd target structure) --- these are exactly the finite $(3\rho)$ terms --- so the factorised form of the logarithms is needed in addition. The explicit cancellation is presented as a confirmation; it is stronger than C-parity, since it also removes an $f^{abc}g_f$ part, which C does not forbid.
\where Section~4.3.

% ======================================================================
\section*{Presentation}

\point{Placeholders (AUTHORS, AFFILIATIONS, ACKNOWLEDGMENTS) and red internal notes on pp.~31, 60--61, 65 and 67.}
\result All red internal notes are removed, and the macro that produced them is deleted. Each note was either resolved (p.~31: collinear cross-check, Major~4; pp.~60--61 and 65: contact power counting, Major~1; p.~67: inclusive constant, now withdrawn) or turned into a plain statement in the new Section~8.1 ``Open points''. The author, affiliation and acknowledgment fields will be completed before submission.

\point{Mixed log/ln; the $\vee$ notation for the cut-off is unusual and hard to read.}
\result The notation is now uniformly $\log$. We have kept $\vee$ for the upper cut-off, for consistency with Lublinsky--Mulian and our previous paper, on which the calculation is built; it is defined at its first use in Section~2.2, $\Lambda<k^+<\vee$, $\vee\equiv\Lambda e^{\delta Y}$, as the upper cut-off on the plus momentum of the soft modes.

\point{No compact final formula \dots\ no plots, no numerical evaluation, and no momentum-space form of the finite terms.}
\result New Section~6.1 gives the complete result in one equation, Eq.~(6.2), with a table that points, for each term (LO$+$RC, BFKL, JIMWLK, DGLAP, the finite terms Row by Row in ordered and colour-singlet form, the odd part, and $\Delta N$), to its equations, together with the groupings that are part of the result and the scheme. It states explicitly that no numerical evaluation and no plot is given. Momentum-space forms are now given for the dilute limit (real emission, Eq.~(7.2); virtual correction for a point-like quark, Eq.~(7.3)) and for $\Delta N$ (Eq.~(F.20)).
\open A numerical evaluation of the finite terms, and a momentum-space form of all finite terms for a general target, are not given (Section~8.1 (iv)).

\point{Dozens of ``verified to $10^{-15}$'' statements cannot be reproduced; the code should be published as ancillary files.}
\result The scripts used for all numerical checks are available from the authors on request; this is now stated in the introduction to Section~7. Where a statement rests on numerics rather than on an analytic argument, the revised text says so.

% ======================================================================
\section*{Required before resubmission}
\begin{itemize}[leftmargin=1.5em]
\item \emph{The dilute-limit comparison with the one-loop Lipatov vertex.} Partly done (Major~2): the real emission agrees exactly with the QMRK cross section, and for a point-like quark the virtual correction has the Regge trajectories and no $\frac{67}{9}$. The comparison of the finite part of the virtual correction with the one-loop vertex, which requires the real emission with matched regulators, is not done and is listed as open.
\item \emph{The $N_f$ (quark-loop) terms, to identify $\beta_0$ unambiguously.} The result remains pure Yang--Mills ($N_f=0$), and the $N_f$ rows are not computed. The identification of $\beta_0$ that the referee asks for was done by the BLM-type device at the level of the kernels (Major~3): all $\frac{11}{3}$, $\frac{11}{6}$, $\frac{11}{12}$ of the virtual and UV terms become $\beta_0$-type; the tail coefficient $-\frac{11}{6}$ does not.
\item \emph{Power counting of all contact kernels.} Done (Major~1): all integrable; the only non-integrable kernel comes with a vanishing colour factor.
\item \emph{A standard collinear cross-check in a stated $\MS$ scheme.} Done (Major~4) for $e^+e^-\to gX$ and $H\to gX$ in CDR with the $\MS$ number-density fragmentation function; $pp\to gX$ not done.
\item \emph{A demonstration that the principal-value, distributional and integration-order prescriptions do not affect the result.} The distributional and integration-order prescriptions do not (Majors~5, 6): the measured-momentum Fourier integrals are unique, the remaining integrals converge absolutely in the given grouping, and the instantaneous terms are independent of the cut-off shape up to a universal shift of the BFKL scale. The principal-value prescription, however, \emph{does} affect the result: the out-state dressing differs by the finite term $\Delta N$, which we now give explicitly, so that both prescriptions are available.
\end{itemize}

We hope that the revised manuscript, which now states precisely what has been established and what remains open, addresses the referee's concerns. We thank the referee again for a report that improved the paper substantially.

\end{document}
